\documentclass[11pt,a4paper]{article}
\usepackage{xcolor}
\usepackage[utf8]{inputenc}
\usepackage{geometry}
\geometry{margin=1in}
\usepackage{hyperref}
\usepackage{enumitem}
\usepackage{amsmath}
\usepackage{listings}

\lstset{
  language=C,
  basicstyle=\ttfamily\small,
  keywordstyle=\bfseries,
  commentstyle=\itshape,
  showstringspaces=false,
  breaklines=true,
  captionpos=b,
  numbers=left,
  numberstyle=\tiny\color{gray},
}

\title{Digital Design Training Program \\ \large Day 8: Embedded Math Algorithms}
\author{Authored by }
\date{April 1, 2026}

\begin{document}

\maketitle

\section*{Theme: Implementing advanced mathematical approximations strictly in C, and exploring extreme bit-level optimizations.}

\subsection*{Module 8.1: RK4 Engine Implementation (60 Minutes)}

\begin{itemize}
    \item \textbf{Goal:} Implement a generic, reusable 4th-Order Runge-Kutta (RK4) engine in C to serve as the core of our custom math library.
    \item \textbf{Conceptual Overview:} Euler's method (Day 7) is simple but accumulates errors quickly. The RK4 method is a significant improvement. Instead of using the slope at the start of an interval, it cleverly samples the slope at four different points within the interval and takes a weighted average. This results in a much more accurate prediction of the next point.
    \item \textbf{Function Pointers:} To make our RK4 engine generic, we need to pass it the derivative function we want to solve (e.g., \verb|y'=y| for \verb|exp(x)|). In C, we can pass functions as arguments using \textbf{function pointers}.
    \begin{itemize}
        \item A normal function declaration is \verb|float my_func(float x, float y);|
        \item A pointer to such a function is declared as \verb|float (*my_func_ptr)(float x, float y);|
        \item We can use \verb|typedef| to create a clean alias: \verb|typedef float (*deriv_func)(float, float);| Now \verb|deriv_func| is a new type representing a pointer to our kind of derivative function.
    \end{itemize}
    \item \textbf{Action Steps \& Code:}
\begin{lstlisting}[caption={Code for Module 8.1: Generic RK4 Solver}]
#include <stdio.h> // For NULL

// Use a typedef to create a new type for our function pointers.
// This type, `deriv_f', is a pointer to any function that
// takes two floats and returns a float.
typedef float (*deriv_f)(float x, float y);

// The core RK4 calculation for a single step.
float rk4_step(float x, float y, float h, deriv_f df) {
    // Calculate the four slopes (k1, k2, k3, k4)
    float k1 = h * df(x, y);
    float k2 = h * df(x + 0.5 * h, y + 0.5 * k1);
    float k3 = h * df(x + 0.5 * h, y + 0.5 * k2);
    float k4 = h * df(x + h, y + k3);
    
    // Return the weighted average
    return y + (k1 + 2.0*k2 + 2.0*k3 + k4) / 6.0;
}

// A wrapper function to solve from a start to an end point.
float rk4_solve(float x0, float y0, float x_target, deriv_f df) {
    float x = x0;
    float y = y0;
    const float h = 0.01; // A reasonably small step size
    int num_steps = (int)((x_target - x0) / h);

    for (int i = 0; i < num_steps; i++) {
        y = rk4_step(x, y, h, df);
        x += h;
    }

    // If x_target was not an exact multiple of h,
    // take one final, smaller step to reach it exactly.
    float remaining_h = x_target - x;
    if (remaining_h > 1e-6) { // Use a small epsilon for float comparison
        y = rk4_step(x, y, remaining_h, df);
    }

    return y;
}

// Example derivative function for y' = y
float deriv_exp(float x, float y) {
    return y;
}

int main(void) {
    // Example of how to use the solver to find e^1
    // float e_approx = rk4_solve(0.0, 1.0, 1.0, deriv_exp);
    // printf(``e^1 is approximately %f\n'', e_approx);
    return 0;
}
\end{lstlisting}
    \item \textbf{Checkpoint:} You have a generic C function, \verb|rk4_solve|, that can numerically approximate the solution to any first-order ODE given a starting point and a target.
\end{itemize}
\newpage
\subsection*{Module 8.2: Trig \& Logs via RK4 and Custom Formatting (60 Minutes)}

\begin{itemize}
    \item \textbf{Goal:} Use the RK4 engine to calculate transcendental functions and develop a lightweight \verb|float|-to-string converter to display the results.
\item \textbf{The Problem with }\verb|sprintf|: The standard C library function \verb|sprintf| is incredibly powerful, but linking in its floating-point conversion module can add over 1.5KB to your program size. On a chip with only 32KB of Flash, this is unacceptable. We must write our own.
    \item \textbf{Action Steps \& Code:}
\begin{lstlisting}[caption={Code for Module 8.2: Custom Float-to-ASCII}]
#include <stdlib.h> // For itoa()

// --- Assume rk4_solve() and lcd driver are included ---

// A lightweight float-to-string function
void ftoa(float val, char* buf, int precision) {
    // Handle negative numbers
    if (val < 0) {
        *buf++ = `-';
        val = -val;
    }

    // Extract the integer part
    long integer_part = (long)val;
    itoa(integer_part, buf, 10); // Convert integer to string

    // Move buffer pointer to the end of the integer string
    while (*buf != `\0') {
        buf++;
    }
    
    // Add the decimal point
    *buf++ = `.';

    // Extract the fractional part
    float fractional_part = val - integer_part;

    // Multiply by 10^precision to make it an integer
    for (int i = 0; i < precision; i++) {
        fractional_part *= 10.0;
    }

    // Convert the new fractional integer to string
    itoa((long)fractional_part, buf, 10);
}

// Derivative for y' = y
float deriv_exp(float x, float y) { return y; }

int main(void) {
    lcd_init();
    lcd_print(``Calculating e^2...'');
    _delay_ms(1000);
    
    // Use the RK4 engine to solve y'=y from y(0)=1 to x=2
    float result = rk4_solve(0.0, 1.0, 2.0, deriv_exp);
    
    char result_str[16];
    ftoa(result, result_str, 3); // Convert with 3 decimal places

    lcd_cmd(0x01); // Clear screen
    lcd_print(``e^2 approx: '');
    lcd_print(result_str);

    while(1) {}
    return 0;
}
\end{lstlisting}
    \item \textbf{Checkpoint:} The LCD displays the numerically integrated value of $e^{2}$ as approximately \verb|7.389|. You have successfully built and displayed the result of a custom math library.
\end{itemize}
\newpage
\subsection*{Module 8.3: The Fast Inverse Square Root (15 Minutes)}

\begin{itemize}
    \item \textbf{Goal:} Witness and implement the legendary Quake III algorithm to see how ``illegal'' bit-level hacking can achieve massive performance gains.
    \item \textbf{Conceptual Overview:} Calculating $1/\sqrt{x}$ is slow. This algorithm, made famous by the game Quake III Arena, achieves it with incredible speed by exploiting the way floating-point numbers are stored in memory. It's a classic example of brilliant, if mind-bending, low-level optimization.
    \item \textbf{The ``Magic'':}
    \begin{enumerate}
        \item It takes the 32-bit float, \verb|y|, and re-interprets its memory location as a 32-bit integer, \verb|i|. This is called pointer casting and is technically ``undefined behavior'' in C, but it works on most compilers.
        \item It performs a single bit-shift (\verb|i >> 1|) and subtracts the result from a ``magic number'' (\verb|0x5f3759df|). This integer operation is a lightning-fast approximation of a logarithm, which effectively calculates a rough inverse square root.
        \item It casts the bits back to a float.
        \item It performs one round of the Newton-Raphson method to polish the initial guess, making it much more accurate.
    \end{enumerate}
    \item \textbf{Action Steps \& Code:}
\begin{lstlisting}[caption={Code for Module 8.3: Fast Inverse Square Root}]
#include <stdint.h> // For int32_t

float Q_rsqrt( float number ) {
    int32_t i;
    float x2, y;
    const float threehalfs = 1.5F;

    x2 = number * 0.5F;
    y  = number;
    
    // Evil floating point bit level hacking
    i  = * ( int32_t * ) &y;
    
    // What the fuck?
    i  = 0x5f3759df - ( i >> 1 );
    
    // Cast bits back to float
    y  = * ( float * ) &i;
    
    // 1st iteration of Newton's method for refinement
    y  = y * ( threehalfs - ( x2 * y * y ) );
    
    return y;
}
\end{lstlisting}
    \item \textbf{Checkpoint:} You can now calculate $1/\sqrt{x}$ almost instantly. This demonstrates the raw power of C pointers and memory casting, showing that sometimes, for performance, you have to break the rules.
\end{itemize}

\end{document}
